\documentclass[10pt,a4paper]{amsart}
\usepackage[margin=19mm]{geometry}
\usepackage{amsmath,amssymb,mathtools}
\usepackage[scr=rsfs,cal=euler]{mathalfa}
\usepackage{xfrac}
\usepackage[unicode,hidelinks,bookmarks=false]{hyperref}
\newcommand{\ie}{i.e.\ }
\newcommand{\cf}{cf.\ }
\newcommand{\resp}{respectively\ }
\newcommand{\Subsection}{\subsection}
\providecommand{\nobreakdash}{\nobreak\hbox{-}}
\setlength{\emergencystretch}{2em}
\multlinegap0pt
\usepackage[shortlabels]{enumitem}
\usepackage{amssymb}
\usepackage{mathtools}
\newtheorem{thm}{Theorem}
\newtheorem{prop}{Proposition}
\newtheorem{rem}{Remark}
\newtheorem{lem}{Lemma}
\newcommand{\ddc}{\mathrm{d} \mathrm{d}^c}
\newcommand{\nat}{\mathbb{N}}
\newcommand{\psh}{\operatorname{PSH}}
\newcommand{\real}{\mathbb{R}}
\newcommand{\sym}{{\rm{Sym}}}
\title{\vspace*{-1.5cm} Kodaira-Iitaka dimension and multiplicity: \\ an analytic perspective}
\author{Siarhei Finski}
\date{Source: arXiv:2603.22194v1 (CC BY 4.0)}
\begin{document}
\maketitle

\noindent\emph{Source and licence.} \vspace*{-1.5cm} Kodaira-Iitaka dimension and multiplicity: \\ an analytic perspective. Version arXiv:2603.22194v1 (CC BY 4.0), distributed by arXiv under the Creative Commons Attribution 4.0 International licence, http://creativecommons.org/licenses/by/4.0/, as recorded in the arXiv licence metadata for this exact version and shown on the abstract page https://arxiv.org/abs/2603.22194v1. Full licence notice: \texttt{licenses/arxiv-2603.22194-CC-BY-4.0.txt}.\par\smallskip

\noindent\textit{Editorial scope.} Counted unit: the article's thm thm\_ki\_form (source0.tex lines 405--412). The article's complete original proof of it is delivered with it (source0.tex lines 1292--1380, one \texttt{proof} environment of 89 lines, which the article places away from the statement). No mathematics is written, repaired or re-derived: the statement and the proof are the article's own lines, in the article's own notation, and no retained line is substituted or re-flowed.  Retained uncounted as the source-local context the proof needs: lines 428--430; lines 437--444; lines 439--443; lines 460--462; lines 977--983; lines 997--1003; lines 1030--1041; lines 1125--1142; lines 1156--1163; lines 1182--1186; lines 1204--1210; lines 1234--1241; lines 1251--1253.  Nothing was omitted from any retained span, so the package carries no omission record.  No retained line contains a citation, so the wrapper carries no bibliography and no bibliography entry.  No retained line contains a figure environment, an image-inclusion command or drawing code, so no image is delivered and the package declares no asset dependency.  Editorial formatting only, in addition: an AMS \texttt{amsart} preamble that restates the article's own declarations and macros used by the retained lines, namely the environments \texttt{thm}, \texttt{prop}, \texttt{rem}, \texttt{lem} on separate counters, together with the standard packages those lines need.  The retained environments print in this document as Theorem 1, Proposition 1, Remark 1, Lemma 1 in order of appearance, whereas the article prints the same items with the numbering of its own full document.  The complete native source of the article as distributed in the arXiv e-print for this exact version is bundled unchanged as \texttt{sources/197004/source0.tex}.
	\begin{equation}\label{eq_dim_ki_yk}
		\kappa(W) = \dim(Y_m), \quad \text{for } m \in \nat \text{ large enough}.
	\end{equation}

	\begin{thm}\label{thm_vol_form}
		For any graded linear series $W \subset R(X, L)$, we have
		\begin{equation}\label{eq_vol_form}
			{\rm{vol}}_{\kappa}(W)
			=
			\frac{1}{\varphi_{*} [\omega]^{n - \kappa(W)}} \int_X c_1(L, P[W, h^L])^{\kappa(W)} \wedge \omega^{n - \kappa(W)}.
		\end{equation}
	\end{thm}

		\begin{equation}
			{\rm{vol}}_{\kappa}(W)
			=
			\frac{1}{\varphi_{*} [\omega]^{n - \kappa(W)}} \int_X c_1(L, P[W, h^L])^{\kappa(W)} \wedge \omega^{n - \kappa(W)}.
		\end{equation}

		\begin{equation}\label{eq_vol_bouck}
			{\rm{vol}}_{\kappa}(L) = \int_X T_{\min}^n.
		\end{equation}

	\begin{thm}\label{thm_wn_monot}
		Let $T_j$, $T'_j$, $j = 1, \ldots, n$, be closed positive $(1, 1)$-currents on $X$ such that $T_j$, $T'_j$ are in the same cohomology class, and $T'_j$ is less singular than $T_j$ for every $j$. 
		Then, we have
		\begin{equation}
			\int_X T_1 \wedge \cdots \wedge T_n \leq \int_X T'_1 \wedge \cdots \wedge T'_n.
		\end{equation}
	\end{thm}

	\begin{thm}\label{thm_cont_begz}
		Let $\alpha_i$, $i \in \{1, \ldots, n\}$ be smooth closed real $(1, 1)$-forms on $X$ representing psef classes. 
		Suppose that $u_i, u_i^k \in \psh(X, \alpha_i)$ are such that $u_i^k$ increase towards $u_j$ outside of a pluripolar subset, as $k \to \infty$.
		Let $\chi_k, \chi : X \to \real$ be quasi-continuous and uniformly bounded such that $\chi_k$ converges to $\chi$, as $k \to \infty$, in capacity.
		We denote $T_i^k := \alpha_i + \ddc u_i^k$ and $T_i := \alpha_i + \ddc u_i$.
		Then $\chi_k \cdot T_1^k \wedge \cdots \wedge T_n^k$ converge weakly towards $\chi \cdot T_1 \wedge \cdots \wedge T_n$, as $k \to \infty$.
	\end{thm}

	\begin{prop}\label{prop_proj_fla}
		Assume the cohomological class of $\alpha$ is big and that $\phi \in \psh(Y, \alpha)$ is bounded.
		Let $\chi$ be a quasi-continuous bounded function on $Y$.
		Then we have
		\begin{equation}
			\int_X \pi^* \chi \cdot (\pi^* \alpha + dd^c \pi^* \phi)^{\kappa} \wedge \omega_0^{n - \kappa}
			=
			\pi_*[\omega_0]^{n - \kappa}
			\cdot
			\int_Y \chi \cdot (\alpha + dd^c \phi)^{\kappa}.
		\end{equation}
	\end{prop}

	\begin{rem}\label{rem_zero_prod}
		a)
		Analogously, we can see that for any $i > \kappa$, we have
		\begin{equation}
			\int_X \pi^* \chi \cdot (\pi^* \alpha + dd^c \pi^* \phi)^{i} \wedge \omega_0^{n - i}
			=
			0.
		\end{equation}
		\par 
		b) The same proof yields that if the $(1, 1)$-forms $\alpha_i$ (resp. $\phi_i \in \psh(Y, \alpha_i)$), $i = 1, \ldots, \kappa$, verify the assumptions of $\alpha$ (resp. $\phi$) from Proposition \ref{prop_proj_fla}, then we have
		\begin{equation}
			\int_X \pi^* \chi \cdot \prod_{i = 1}^{\kappa} (\pi^* \alpha_i + dd^c \pi^* \phi_i) \wedge \omega_0^{n - \kappa}
			=
			\pi_*[\omega_0]^{n - \kappa}
			\cdot
			\int_Y \chi \cdot  \prod_{i = 1}^{\kappa} (\alpha_i + dd^c \phi_i).
		\end{equation}
	\end{rem}

	\begin{thm}\label{thm_sym_subalgebra}
		For any $\epsilon > 0$, there is $m_0 \in \nat$ such that for any $m \geq m_0$, there is $k_0 \in \nat$ such that for any $k \geq k_0$, we have
		\begin{equation}
			\frac{\dim [\sym^k W_m]}{\dim W_{km}}
			\geq 
			1 - \epsilon.
		\end{equation}
	\end{thm}

	\begin{equation}\label{eq_sym_yk_rel}
		[\sym^k W_m]
		\simeq
		H^0(Y_m, A_m^{\otimes k}).
	\end{equation}

	\begin{equation}\label{eq_incl_resol_coh}
		H^0(Y_m, A_m^{\otimes k}) 
		\overset{p_m^*}{\hookrightarrow} 
		H^0(\widehat{Y}_m, \widehat{A}_m^{\otimes k}) 
		\overset{\widehat{\varphi}_m^*}{\hookrightarrow} 
		H^0(\widehat{X}_m, \widehat{L}^{\otimes k m}).
	\end{equation}

	\begin{equation}\label{eq_fs_n_k_a_k_rel}
		\pi_m^* FS(N_m)
		=
		\varphi_m^*
		\widehat{h}^A_m
		\cdot
		h^{E_m}_{{\rm{sing}}}.
	\end{equation}

	\begin{lem}\label{lem_fs_sm}
		The sequence of the Fubini-Study (singular) metrics $FS(N_k)$, $k \in \nat^*$, is submultiplicative for any submultiplicative graded norm $N = (N_k)_{k = 0}^{\infty}$, i.e. for any $k, l \in \nat$, we have $FS(N_{k + l}) \leq FS(N_k) \cdot FS(N_l)$.
	\end{lem}

	\begin{thm}\label{thm_ki_form}
		For any graded linear series $W \subset R(X, L)$, we have
		\begin{equation}\label{eq_kappa_w}	
			\kappa(W) = \max\Big\{
				i \in \{ 0, 1, \ldots, n \} :  c_1(L, P[W, h^L])^{i} \wedge \omega^{n - i} \neq 0
			\Big\}.
		\end{equation}
	\end{thm}

	\begin{proof}[Proof of Theorems \ref{thm_ki_form} and \ref{thm_vol_form}]
		First of all, by Lemma \ref{lem_fs_sm}, the sequence of singular metrics
		\begin{equation}
			h^L_{2^k}
			:=
			FS({\textrm{Ban}}_{2^k}^{\infty}[W](h^L))^{\frac{1}{2^k}}
		\end{equation}
		is non-increasing.
		From now on, for brevity, we denote $m := 2^k$.
		By Theorem \ref{thm_cont_begz}, for any $i \in \nat$, we obtain the following convergence
		\begin{equation}
			\int_X c_1(L, P[W, h^L])^{i} \wedge \omega^{n - i}
			=
			\lim_{m \to \infty} \int_X c_1(L, h^L_m)^{i} \wedge \omega^{n - i}.
		\end{equation}
		We denote by $\widehat{h}^A_m$ the metric on $\widehat{A}_{m}$ induced by ${\textrm{Ban}}_{m}^{\infty}[W](h^L)$ as described before (\ref{eq_fs_n_k_a_k_rel}).
		By (\ref{eq_fs_n_k_a_k_rel}), we have
		\begin{equation}
			m \cdot \pi_m^* c_1(L, h^L_m)
			=
			\widehat{\varphi}_{m}^* c_1(\widehat{A}_{m}, \widehat{h}^A_m)
			+
			[E_m],
		\end{equation}
		where $[E_m]$ is the current of integration along $E_m$.
		From this and the fact that the non-pluripolar product doesn't charge the analytic subsets (as $E_m$), we deduce 
		\begin{equation}\label{eq_form_pullback}
			m^i \cdot \pi_m^* c_1(L, h^L_m)^{i} \wedge \pi_m^* \omega^{n - i}
			=
			\widehat{\varphi}_m^* c_1(\widehat{A}_m, \widehat{h}^A_m)^{i} \wedge \pi_m^* \omega^{n - i}.
		\end{equation}
		Recall that by (\ref{eq_dim_ki_yk}), we have $\dim \widehat{Y}_m = \kappa(W)$, for $m$ large enough.
		From this, Remark \ref{rem_zero_prod} and (\ref{eq_form_pullback}), we deduce immediately that for $i > \kappa(W)$, we have
		\begin{equation}\label{eq_i_big_zero}
			\int_X c_1(L, h^L_m)^{i} \wedge \omega^{n - i}
			=
			0.
		\end{equation}
		\par 
		Moreover, by Proposition \ref{prop_proj_fla} and (\ref{eq_form_pullback}), we have
		\begin{equation}\label{eq_int_l_a_id}
			m^{\kappa(W)} \cdot \int_X c_1(L, h^L_m)^{\kappa(W)} \wedge \omega^{n - \kappa(W)}
			=
			\widehat{\varphi}_{*}[ \pi_m^* \omega]^{n - \kappa(W)}
			\cdot
			\int_{\widehat{Y}_m} c_1(\widehat{A}_m, \widehat{h}^A_m)^{\kappa(W)},
		\end{equation}
		where $\widehat{\varphi}_{*}[ \pi_m^* \omega]^{n - \kappa(W)}$ is defined analogously to the quantity appearing in (\ref{eq_vol_form}).
		\par 
		Since $p_m$ and $\pi_m$ are birational modifications, we deduce that for $m$ large enough
		\begin{equation}\label{eq_bir_fib_integg}
			\widehat{\varphi}_{*}[ \pi_m^* \omega]^{n - \kappa(W)}
			=
			\varphi_{*} [\omega]^{n - \kappa(W)}.
		\end{equation}
		\par 
		Also, since $\widehat{h}^A_m$ is a continuous metric with a psh potential on $\widehat{A}_m$, we have
		\begin{equation}\label{eq_y_m_int_a_vol}
			\int_{\widehat{Y}_m} c_1(\widehat{A}_m, \widehat{h}^A_m)^{\kappa(W)}
			=
			{\rm{vol}}(\widehat{A}_m),
		\end{equation}
		by a version of (\ref{eq_vol_bouck}).
		Note, however, that in the current semi-ample setting the corresponding result follows immediately from Theorem \ref{thm_wn_monot} and the Hilbert-Samuel theorem.
		\par 
		Also, by Theorem \ref{thm_sym_subalgebra}, (\ref{eq_sym_yk_rel}) and (\ref{eq_incl_resol_coh}), we deduce that for any $\epsilon > 0$, there is $m_0 \in \nat$, so that for any $m \geq m_0$, we have
		\begin{equation}\label{eq_bol_w_a_est}
			(1 + \epsilon) 
			\cdot 
			m^{\kappa(W)}
			\cdot
			{\rm{vol}}_{\kappa}(W)
			\geq
			{\rm{vol}}(\widehat{A}_m)
			\geq
			(1 - \epsilon) 
			\cdot 
			m^{\kappa(W)}
			\cdot
			{\rm{vol}}_{\kappa}(W).
		\end{equation}
		By combining (\ref{eq_form_pullback}), (\ref{eq_int_l_a_id}), (\ref{eq_bir_fib_integg}), (\ref{eq_y_m_int_a_vol}) and (\ref{eq_bol_w_a_est}), we establish immediately Theorem \ref{thm_vol_form}.
		\par 
		Since ${\rm{vol}}_{\kappa}(W) > 0$, by Theorem \ref{thm_vol_form}, we deduce that
		\begin{equation}
			\int_X c_1(L, P[W, h^L])^{\kappa(W)} \wedge \omega^{n - \kappa(W)} > 0.
		\end{equation}
		This along with (\ref{eq_i_big_zero}) finishes the proof of Theorem \ref{thm_ki_form}.
	\end{proof}

\end{document}
